\section{The actual preferred-cusp
translation}\label{pl:cusp}

We use the actual parabolic-main face maps in \cite[Section 7]{StrictGeometry}.
They also compute the cusp tangent of the all-circle family: at $M=1$
the native chart of \cite{Nonparabolic} has its preliminary white
coordinate equal to $X$, and its four white matrices are the matrices
below up to the stated projective normalizations. The full ordered disk
conjugacies identify the shaded transports as well. Alternatively, the
positive marked gluing germ identifies both constructions with the same
original complete holonomy for prescribed auxiliary longitude logs.
Both descriptions retain the same physical wire labels. Thus the
following calculation is of the actual main peripheral word, not a
comparison of critical values.

\subsection{Meridians, endpoint signs, and additive
displacement}\label{meridians-endpoint-signs-and-additive-displacement}

Normalize each ordered quadrilateral to
\((C_0,C_1,C_2,C_3)=(0,\infty,1,R)\), and put \(E=e^{2\pi i b}\). In
these charts, a white face is a piece of an original projection region;
the other, shaded faces are pieces of the crossing disks. Use the actual
white matrices \begin{equation}
\begin{aligned}
W_0&=\begin{pmatrix}1&0\\0&E\end{pmatrix},
&W_1&=\begin{pmatrix}1&E-1\\0&E\end{pmatrix},\\
W_2&=\begin{pmatrix}E-R&R(1-E)\\E-1&1-ER\end{pmatrix},
&W_3&=\begin{pmatrix}R&0\\1-E&ER\end{pmatrix}.
\end{aligned}
\label{plc:4}
\end{equation} At an open shaded boundary some of these are completion maps. The
marked geometry input in Proposition~\ref{pl:imports} proves that every occurrence used in the final
closed peripheral paths is supplied by the actual adjacent white face.

Set \(d_j=(-1)^j\) and \begin{equation}
\mu_j=(W_{j-1}^{-1}W_j)^{d_j}.
\label{plc:5}
\end{equation} The coordinates \begin{equation}
\begin{aligned}
\xi_0(z)&=\frac{ER}{(E-1)z},
&\xi_1(z)&=-\frac{z}{E-1},\\
\xi_2(z)&=-\frac{E(R-1)}{(E-1)(z-1)},
&\xi_3(z)&=-\frac{R(R-1)}{(E-1)(z-R)}
\end{aligned}
\label{plc:6}
\end{equation} send the corresponding corner to infinity and send \(\mu_j\) to
\(u\mapsto u+1\). These identities follow by multiplying \eqref{plc:4}.

The endpoint sign has a geometric meaning. Orient a main arc into its
leaf at corner \(j\). Its positively oriented small meridian is
\(\mu_j^{d_j}\); orient the arc out of the leaf and it is
\(\mu_j^{-d_j}\). This is the right-hand meridian convention in the
ordered projection chart \(0,1,2,3\): crossing one white sector and
returning through the preceding one gives \(W_{j-1}^{-1}W_j\) for an
inward arc. Formula \eqref{plc:5} therefore supplies exactly the displayed signs.
At infinity use the finite local coordinate \(1/z\). Reversing the
strand orientation reverses both its meridian and the direction of its
longitude.

Consider an open main arc from corner \(j\) to corner \(k\), including
its specified peripheral push-off. In the local coordinates its actual
transport has the form \begin{equation}
u_k=-d_jd_k\,u_j+\beta.
\label{plc:7}
\end{equation} The multiplier follows by transporting the same oriented meridian: it
is \(d_j\) times the local unit at the incoming endpoint and \(-d_k\)
times the local unit at the outgoing endpoint. Define the displacement
of this arc to be \begin{equation}
D_{\mathrm{arc}}=-d_k\beta.
\label{plc:8}
\end{equation} This definition does not depend on the direction in which the arc is
followed. The inverse affine map has translation \(-\beta/(-d_jd_k)\),
and its displacement is \(-d_j[-\beta/(-d_jd_k)]=-d_k\beta\). This
simultaneously reverses the path and the oriented meridian; it does not
add raw translations from incompatible meridian frames.

When two endpoints are joined by an actual identity segment, its
coordinate change from \(j\) to \(k\) has the same affine multiplier
\(-d_jd_k\). Its contribution is instead \begin{equation}
D_{\mathrm{join}}=d_k\beta.
\label{plc:9}
\end{equation} Here the path leaves the first arc and enters the next. The same
inverse-map calculation proves direction independence.

These displacements add under concatenation. For an arc \(i\to j\), a
joining segment \(j\to k\), and an arc \(k\to l\), write their shifts as
\(\beta_X,\beta_J,\beta_Y\). Multiplying the three affine maps gives
shift \begin{equation}
(-d_kd_l)(-d_jd_k)\beta_X
       +(-d_kd_l)\beta_J+\beta_Y.
\end{equation} Multiplying by \(-d_l\) gives exactly \begin{equation}
-d_j\beta_X+d_k\beta_J-d_l\beta_Y.
\label{plc:10}
\end{equation} The product multiplier is \(-d_id_l\), as required. The same identity
applies when a join creates a closed component. A closed component has
multiplier one and its displacement is its actual longitude translation
in the consistently oriented unit-meridian frame.

Let \(S(Q)\) be the sum of these displacements for the two open arcs and
any already closed main components inside a tangle. If a boundary
coordinate at a corner changes by \(u'_j=\alpha_j u_j+\beta_j\), with
\(\alpha_j=d_j/d'_j\), then \begin{equation}
S(Q')-S(Q)=-\sum_{\text{four endpoints }j}d'_j\beta_j.
\label{plc:11}
\end{equation} This follows from \eqref{plc:8}, once at each endpoint. Closed components do
not change. Equations \eqref{plc:7}--\eqref{plc:11} are a local chain lemma: they apply
regardless of how many odd leaves a component visits or which direction
is used to evaluate it.

For an open expression, this is boundary-chart bookkeeping in the
original two-ball cut, with its completion maps still labeled as such.
It does not declare its metric quadrilateral to be a physical
four-punctured tangle boundary or assign virtual white maps to new group
generators. The product identities transport this bookkeeping. At the
final closure every used white completion is the actual adjacent
original white transition, by the marked geometry input. Only then do the
assembled paths give the asserted closed peripheral words.

Equivalently, for any one step write
\(\xi_bF\xi_a^{-1}(w)=\kappa w+\beta\), and let the globally oriented
meridian be the local unit raised to \(\alpha_a,\alpha_b\in\{1,-1\}\).
Consistency is \(\alpha_b\kappa=\alpha_a\); the actual translation
contribution is \(\alpha_b\beta\). Reversing the traversal replaces the
shift by \(-\beta/\kappa\), and reversing the oriented meridian replaces
the target sign by \(-\alpha_a\). Their product is again
\(\alpha_b\beta\). Equations \eqref{plc:8} and \eqref{plc:9} are this rule at an arc and a
joining segment respectively. This explicitly includes the
multiplier-minus-one odd routes below.

\subsection{The exact product
potential}\label{the-exact-product-potential}

Define the boundary potential \begin{equation}
P(E,R)=\frac{E+R}{E-1},\qquad I(Q)=S(Q)-P(E,R).
\label{plc:12}
\end{equation} For the original ordered product prism we claim \begin{equation}
I(FG)=I(F)+I(G).
\label{plc:13}
\end{equation} This is an identity of actual coordinate transports, not an equality
of their absolute values.

Here is a rational verification retaining every translation. Write the
prism angle phases as \(a,b,c\), so its vertices are \begin{equation}
v_0=1,\quad v_1=a,\quad v_2=ab,\quad v_3=abc,\qquad
 P_*=
 1+a-a\,\frac{ab(1+c)-1-a}{a^2b^2c-a}.
\label{plc:14}
\end{equation} The lower, upper and output ordered quadrilaterals are
\((1,a,P_*,\infty)\), \((\infty,P_*,ab,abc)\), and \((1,a,ab,abc)\).
Their \(E\) parameters are \(bc,ab,abc\). Their normalized fourth
vertices are \begin{equation}
\begin{aligned}
 R_F&=\frac{a(b-1)(bc-1)}{(ab-1)(abc-1)},\\
 R_G&=\frac{c(b-1)(ab-1)}{(bc-1)(abc-1)},\\
 R_o&=\frac{(b-1)(abc-1)}{(ab-1)(bc-1)}.
\end{aligned}
\label{plc:15}
\end{equation} Apply \eqref{plc:6} after the unique normalized three-point map for each
quadrilateral. The two joining and four external rebasing maps all have
multiplier one. Their translations are \begin{equation}
\begin{array}{c|c}
F_2\longrightarrow G_1&bc/(bc-1)\\
F_3\longrightarrow G_0&-a(b-1)/[(ab-1)(abc-1)]\\
F_0\longrightarrow o_0&0\\
F_1\longrightarrow o_1&0\\
G_2\longrightarrow o_2&ab(c-1)/[(ab-1)(abc-1)]\\
G_3\longrightarrow o_3&(b-1)(c-1)/[(ab-1)(bc-1)(abc-1)].
\end{array}
\label{plc:16}
\end{equation} For example a world quadrilateral \((q_0,q_1,q_2,q_3)\) is normalized
by \((z-q_0)(q_2-q_1)/[(z-q_1)(q_2-q_0)]\), with the corresponding
limits at infinity. Inserting this map into \eqref{plc:6} gives \eqref{plc:16}.

By \eqref{plc:9} and \eqref{plc:11}, the correction \(S(FG)-S(F)-S(G)\) is minus the first
entry of \eqref{plc:16}, plus the second, minus the fifth and plus the sixth.
Substituting \eqref{plc:15} into \eqref{plc:12} gives exactly that expression: \begin{equation}
S(FG)-S(F)-S(G)
                  =P(abc,R_o)-P(bc,R_F)-P(ab,R_G).
\label{plc:17}
\end{equation} This proves \eqref{plc:13}. All variables in \eqref{plc:14}--\eqref{plc:17} are independent
wherever the displayed charts are nondegenerate. Thus the identity
continues to the original complex deformed branch. In particular hidden
binary resolutions inside a product hub contribute no uncancelled term.

\subsection{Turns and the final
closure}\label{turns-and-the-final-closure}

The actual quarter-turn is \begin{equation}
T_Q(z)=\frac{R-1}{z-1},\qquad
 (E,R)\longmapsto(E^{-1},1-R),\qquad j\longmapsto j-1.
\label{plc:18}
\end{equation} In the coordinates \eqref{plc:6}, the four rebasing maps have multiplier minus
one. In old corner order \(j=0,1,2,3\), their shifts are \begin{equation}
\frac{ER}{E-1},\quad-\frac1{E-1},\quad0,\quad\frac R{E-1}.
\label{plc:19}
\end{equation} Equation \eqref{plc:11} gives \(S(TQ)-S(Q)=(ER-R+1)/(E-1)\). Subtract the two
potentials in \eqref{plc:12}. The result is \begin{equation}
I(TQ)-I(Q)=\frac{2(E+1)}{E-1}=-2i\cot(\pi b).
\label{plc:20}
\end{equation} This calculation retains the minus multipliers in \eqref{plc:19}; summing their
shifts without \eqref{plc:11} would give the wrong sign.

At the final reverse-order closing gluing, the two input parameters are
\((E,R)\) and \((E^{-1},R)\), and the actual map from the second to the
first is \begin{equation}
C_Q(z)=\frac{z-R}{z-1},\qquad j\longmapsto3-j.
\label{plc:21}
\end{equation} Its four affine multipliers are one. In old corner order their shifts
are \begin{equation}
\frac R{E-1},\quad-\frac E{E-1},
                       \quad-\frac1{E-1},\quad\frac{ER}{E-1}.
\label{plc:22}
\end{equation} The joining correction \eqref{plc:9} is \(R-1\), whereas
\(P(E,R)+P(E^{-1},R)=1-R\). Therefore the total translation of the
closed main components is exactly \begin{equation}
I(F)+I(G).
\label{plc:23}
\end{equation} The original closed component decomposition is retained throughout.
When the main strands form a knot, \eqref{plc:23} is the translation of one
primitive longitudinal traversal.

\subsection{Pure and odd leaf paths}\label{pure-and-odd-leaf-paths}

The actual pure and odd shaded maps are \begin{equation}
S=\begin{pmatrix}0&-ER\\1&-1-E\end{pmatrix},\quad
 A=\begin{pmatrix}R&-R\\1&ER-1-E\end{pmatrix},\quad
 B=\begin{pmatrix}R/E-1-1/E&R\\-1&R\end{pmatrix}.
\label{plc:24}
\end{equation} The pure arcs \(0\to1\) and \(2\to3\) use \(S^{-1}\) and \(S\). From
\eqref{plc:6}, the translations of \(S:1\to0\) and \(S:2\to3\) are \((E+1)/(E-1)\)
and \((R-1)/(E-1)\). Thus \begin{equation}
S(h_0)=\frac{R-E-2}{E-1},\qquad
 I(h_0)=-\frac{2(E+1)}{E-1}=2i\cot(\pi x).
\label{plc:25}
\end{equation}

At an odd leaf the main arcs pair \(0\) with \(2\) and \(1\) with \(3\).
The shaded maps \(B:0^+\to2^-\), \(A:1^+\to3^-\) switch projection
balls. To express a path again in the plus frame, use the indicated
actual adjacent white return. For the physical sign \(\epsilon=+1\),
choose \begin{equation}
L_{02}^{+}=W_1^{-1}B,\qquad L_{13}^{+}=W_2^{-1}A.
\label{plc:26}
\end{equation} For \(\epsilon=-1\), choose \begin{equation}
L_{02}^{-}=W_2^{-1}B,\qquad L_{13}^{-}=W_3^{-1}A.
\label{plc:27}
\end{equation} These are specified peripheral paths. We will determine their framing
by homology rather than assuming that an adjacent white-side choice is
the preferred longitude.

In the plus \(\xi\) coordinates, all four maps have multiplier minus
one, as \eqref{plc:7} requires for equal-parity endpoints. Their shifts, in the
order \(02,13\), are \begin{equation}
\begin{aligned}
\epsilon=+1:\quad&
  \frac E{E-1},\quad\frac{R-E-1}{E-1},\\
\epsilon=-1:\quad&
  \frac1{E-1},\quad\frac{R-2}{E-1}.
\end{aligned}
\label{plc:28}
\end{equation} Using \eqref{plc:8}, rather than summing these raw shifts, gives \begin{equation}
S(h_\epsilon)=S(h_0)-\epsilon,\qquad
                  I(h_\epsilon)=2i\cot(\pi x)-\epsilon.
\label{plc:29}
\end{equation} The relations hold at independent nearby \(E,R\). The selected hidden
auxiliary color affects \(R\) but does not change \eqref{plc:29}.

Every original \(h_\epsilon\) is the turn of its unturned
crossing-circle piece \(p_\epsilon\). Its input label is \(1-x\).
Equation \eqref{plc:20} therefore says \(I(p_\epsilon)=-\epsilon\). Equations
\eqref{plc:13}, \eqref{plc:20} and \eqref{plc:23} now prove for the specified closed longitudinal path
\begin{equation}
\tau_{\mathrm{path}}
        =2i\sum_{e\in\mathcal E_T}\cot(\pi x_e)-\sum_i\epsilon_i.
\label{plc:30}
\end{equation} This explains precisely why the sum is over physical turn edges, not
arbitrary unturned edges in a binary product resolution. Lowering a turn
makes its covariant index \(x_e\) complementary to its unturned input.
Consecutive double turns contribute opposite cotangents and may be
suppressed. Root changes merely reassociate this same geometric sum.

\subsection{The integer framing of the constructed
path}\label{the-integer-framing-of-the-constructed-path}

The original marked threading calculation identifies the actual
auxiliary meridians as \begin{equation}
M_+=W_0^{-1}B\quad\text{for physical }R,\qquad
 M_-=A^{-1}W_0\quad\text{for physical }R^{-1}.
\label{plc:31}
\end{equation} These signs are fixed geometrically in the marked geometric
deformation \cite[Sections 5 and 8]{StrictGeometry}, from the original threading word; they are not
chosen from a quantum asymptotic.

The exact odd edge relations give \begin{equation}
\begin{array}{c|cc}
 &L_{02}&L_{13}\\ \hline
R&\mu_1M_+&M_+^{-1}\\
R^{-1}&M_-&\mu_0M_-^{-1}.
\end{array}
\label{plc:32}
\end{equation} For example \(W_2^{-1}A=M_+^{-1}\), and
\(W_1^{-1}B=(W_1^{-1}W_0)M_+=\mu_1M_+\). For the second row use
\(W_2^{-1}BW_0^{-1}A=I\) and \(W_3^{-1}W_0=\mu_0\). These are identities
of actual group paths, using the already established edge words, and
hence also identities in homology.

Orient the main knot. At an odd leaf let \(t_0,t_1\in\{1,-1\}\) record
whether it follows the canonical arcs \(0\to2\), \(1\to3\) or their
reverses. By the endpoint convention in Section~\ref{pl:cusp}, \begin{equation}
[\mu_0]=t_0[\mu_K],\qquad
                        [\mu_1]=-t_1[\mu_K].
\label{plc:33}
\end{equation} The first row of \eqref{plc:32}, traversed with these signs, contributes \begin{equation}
-t_0t_1[\mu_K]+(t_0-t_1)[M_+]
\label{plc:34}
\end{equation} to the longitudinal homology. The second contributes \begin{equation}
t_0t_1[\mu_K]+(t_0-t_1)[M_-].
\label{plc:35}
\end{equation} The same coefficient \(t_0-t_1\) is the oriented intersection of the
main strands with the original crossing disk. Thus it is their linking
number with the circle.

The main-meridian coefficients in \eqref{plc:34}--\eqref{plc:35} are exactly the oriented
writhe contributions. To check the convention, undo the literal
quarter-turn \(h=T(R^\epsilon)\). Parallel upward strands of the
unturned positive braid become \(t_0=1,t_1=-1\): one diagonal is
reversed by this quarter-turn. Hence the positive crossing has sign
\(-t_0t_1\), and the negative crossing has sign \(t_0t_1\). Reversing
just one strand reverses both signs. This is the usual oriented crossing
definition, applied to the actual ordered terminal map.

For a pure leaf its auxiliary meridian is \(M_0=S^{-1}\). Its two paths
are \(M_0,M_0^{-1}\), so their homology contributes only
\((t_0-t_1)[M_0]\), and no main-meridian term. Products and turns
outside the leaf boxes are actual crossing-free identity strips and
rotations. They do not introduce an additional homology contribution.
Thus the full path from \eqref{plc:30} has \begin{equation}
[L_{\mathrm{path}}]
          =w_\epsilon[\mu_K]+\sum_i\ell_i[\mu_i].
\label{plc:36}
\end{equation} It traverses the main knot once because every original segment and
disk puncture is followed once in its actual oriented component.
Therefore it differs from the preferred augmented longitude by exactly
\(w_\epsilon\) main meridians: \begin{equation}
\lambda_K=L_{\mathrm{path}}\mu_K^{-w_\epsilon},\qquad
        \tau_{\lambda_K}
          =2i\sum_e\cot(\pi x_e)-\sum_i\epsilon_i-w_\epsilon.
\label{plc:37}
\end{equation} This step determines the integer part; parabolicity alone would not
do so.
